\dissabstract

Meniscal injuries are among the most frequently encountered knee pathologies in orthopedic practice, often resulting in chronic pain, joint dysfunction, and long-term degenerative changes. Despite the meniscus's essential role in absorbing axial loads and maintaining joint stability, current treatment strategies—including partial meniscectomy, suture repair, and synthetic implants—fail to replicate the native tissue's mechanical function and viscoelastic behavior. These limitations frequently lead to altered joint kinematics, cartilage degeneration, and early-onset osteoarthritis. There remains a critical unmet need for a scalable, cost-effective, and mechanically robust meniscal replacement that also supports biological integration.

This dissertation presents a novel, clinically relevant scaffold fabrication~approach- Hybrid Hydrogels Augmented via Additive Network Integration (HANI)—\\which enables the development of composite scaffolds with customizable biomechanical profiles. The HANI platform leverages gelatin-based hydrogels, crosslinked with glutaraldehyde (GTA), and augmented with poly-$\varepsilon$-caprolactone (PCL) reinforcement structures fabricated via cost-effective additive manufacturing techniques, including Fused Deposition Modeling (FDM) and Stereolithography (SLA). Custom-fabricated molds and thermal alignment masks were developed to guide the seamless~integration of anisotropic PCL architectures—including aligned, circumferential, and 3D thermo\-molded grids—into the hydrogel matrix.

A modular testing apparatus was concurrently developed to evaluate scaffold mechanical performance under confined compression. This system enabled high-resolution characterization of stress-relaxation behavior, peak stress response, and equilibrium moduli across variable cross-linking densities and PCL configurations. Experimental results demonstrated that HANI scaffolds could be fine-tuned to match the poro\-viscoelastic properties of the native meniscus, with reinforced constructs showing significantly improved compressive stiffness, stress dissipation, and~structural anisotropy.

The clinical impact of this work lies in the scalability and translational readiness of the HANI methodology. By utilizing low-cost materials and widely accessible fabrication technologies, the approach supports rapid prototyping of patient-specific implants without the complexity or expense of bioprinting. Furthermore, the modularity of the HANI platform allows for adaptation across a spectrum of orthopedic soft tissue engineering applications—including cartilage, intervertebral disc, and tendon constructs—thereby broadening its therapeutic relevance.

Altogether, this research introduces a transformative scaffold design paradigm that harmonizes mechanical fidelity with biological compatibility, offering a path forward for next-generation, tissue-engineered solutions that are affordable, adaptable, and clinically actionable.
    
    